
\documentclass{article}
\usepackage{colortbl}
\usepackage{makecell}
\usepackage{multirow}
\usepackage{supertabular}

\begin{document}

\newcounter{utterance}

\twocolumn

{ \footnotesize  \setcounter{utterance}{1}
\setlength{\tabcolsep}{0pt}
\begin{supertabular}{c@{$\;$}|p{.15\linewidth}@{}p{.15\linewidth}p{.15\linewidth}p{.15\linewidth}p{.15\linewidth}p{.15\linewidth}}

    \# & $\;$A & \multicolumn{4}{c}{Game Master} & $\;\:$B\\
    \hline 

    \theutterance \stepcounter{utterance}  

    & & \multicolumn{4}{p{0.6\linewidth}}{\cellcolor[rgb]{0.9,0.9,0.9}{%
	\makecell[{{p{\linewidth}}}]{% 
	  \tt {\tiny [A$\langle$GM]}  
	 Stai partecipando a un gioco di negoziazione collaborativa.\\ \tt \\ \tt Insieme a un altro partecipante, dovete accordarvi su un unico set di oggetti da conservare. Ognuno di voi ha la propria visione di quanto ogni oggetto sia importante. Non sai come l'altro partecipante valuta gli oggetti. Inoltre, ti viene fornito lo sforzo che ogni oggetto richiede.\\ \tt Potete accordarvi su un set solo se lo sforzo totale degli oggetti selezionati non supera un limite condiviso:\\ \tt \\ \tt LIMITE = 3288\\ \tt \\ \tt Ecco i valori di sforzo dei singoli oggetti:\\ \tt \\ \tt Sforzo oggetto = {"C08": 329, "C55": 41, "A43": 417, "A67": 75, "B42": 390, "A62": 887, "B19": 808, "C43": 151, "C73": 849, "C10": 129, "B36": 350, "A18": 118, "C61": 630, "C67": 602, "C42": 801}\\ \tt \\ \tt Ecco la tua visione personale sull'importanza di ogni oggetto:\\ \tt \\ \tt Valori di importanza dell'oggetto = {"C08": 138, "C55": 583, "A43": 868, "A67": 822, "B42": 783, "A62": 65, "B19": 262, "C43": 121, "C73": 508, "C10": 780, "B36": 461, "A18": 484, "C61": 668, "C67": 389, "C42": 808}\\ \tt \\ \tt Obiettivo:\\ \tt \\ \tt Il tuo obiettivo è negoziare un set condiviso di oggetti che ti avvantaggi il più possibile (cioè, massimizzi l'importanza totale per TE), rimanendo entro il LIMITEE. Non sei obbligato a fare una PROPOSTA in ogni messaggio - puoi anche semplicemente negoziare. Tutte le tattiche sono permesse!\\ \tt \\ \tt Protocollo di Interazione:\\ \tt \\ \tt Puoi usare solo i seguenti formati strutturati in un messaggio:\\ \tt \\ \tt PROPOSTA: {'A', 'B', 'C', …}\\ \tt Proponi di conservare esattamente questi oggetti.\\ \tt RIFIUTO: {'A', 'B', 'C', …}\\ \tt Rifiuta esplicitamente la proposta dell'avversario.\\ \tt ARGOMENTAZIONE: {'...'}\\ \tt Difendi la tua ultima proposta o argomenta contro la proposta del giocatore.\\ \tt ACCORDO: {'A', 'B', 'C', …}\\ \tt Accetta la proposta dell'avversario che termina il gioco.\\ \tt RAGIONAMENTO STRATEGICO: {'...'}\\ \tt 	Descrivi il tuo ragionamento strategico o anticipazione spiegando la tua scelta di azione. Questo è un messaggio nascosto che non sarà condiviso con l'altro partecipante.\\ \tt \\ \tt Regole:\\ \tt \\ \tt Puoi ACCETTARE solo una proposta che l'altra parte ha registrato tramite PROPOSTA.\\ \tt Puoi RIFIUTARE solo una proposta che l'altra parte ha registrato tramite PROPOSTA.\\ \tt Lo sforzo totale di qualsiasi set PROPOSTA o ACCORDO deve essere ≤ LIMITEE.\\ \tt NON rivelare i tuoi punteggi di importanza nascosti.\\ \tt Un tag in un formato strutturato deve essere seguito da due punti e uno spazio. L'argomento deve essere un set python contenente 0 o più stringhe.\\ \tt Quindi, deve essere della forma TAG: {...}\\ \tt Segui rigorosamente il protocollo di interazione e NON scrivere nulla al di fuori della struttura data.\\ \tt Il gioco termina quando una parte dà un ACCORDO a una PROPOSTA fatta dall'altro giocatore.\\ \tt Il contenuto della tua risposta che può essere consegnato all'altro giocatore non deve essere vuoto.\\ \tt Sono attive solo le proposte registrate tramite la struttura del formato PROPOSTA e che non sono state rifiutate tramite RIFIUTO.\\ \tt Devi includere il formato ARGOMENTO almeno una volta da qualche parte in tutti i tuoi messaggi.\\ \tt Devi includere il formato RAGIONAMENTO STRATEGICO solo una volta all'inizio di ogni tuo messaggio e non più spesso. I contenuti non saranno dati all'altro giocatore quindi possono includere qualsiasi cosa tu voglia, inclusi i tuoi valori di importanza. Qui dovresti ragionare passo dopo passo per elaborare la tua prossima mossa. 
	  } 
	   } 
	   } 
	 & \\ 
 

    \theutterance \stepcounter{utterance}  

    & \multicolumn{4}{p{0.6\linewidth}}{\cellcolor[rgb]{0.8,1,0.9}{%
	\makecell[{{p{\linewidth}}}]{% 
	  \tt {\tiny [A$\rangle$GM]}  
	 RAGIONAMENTO STRATEGICO: {'Calcolo pesi valore/sforzo per massimizzare la mia utilità entro 3288. Rapporto valore/effort (circa): C55 14.2, A67 11.0, C10 6.0, A43 2.08, C61 1.06, C42 1.01, A18 4.1, C08 0.42, B42 2.01, B36 1.32, C43 0.80, C67 0.65, C73 0.60, B19 0.32, A62 0.07. Un greedy per rapporto dà top: C55, A67, C10, A18, A43, poi C61/C42, B42, B36, ecc. Controllo combinazioni: Un set forte è {C55(41), A67(75), C10(129), A18(118), A43(417), B42(390), B36(350), C61(630), C42(801)} totale effort 3051 e valore 583+822+780+484+868+783+461+668+808=6257. Rimane margine 237: posso aggiungere C43(151) valore 121 -> totale 3202 valore 6378. Rimangono 86, nessun altro aggiungibile (minimo 329). Potrei provare anche C67(602) al posto di qualcosa, ma peggiora. Un altro set: sostituire B36 con C67? peggiora. Quindi proposta iniziale: {C55, A67, C10, A18, A43, B42, C61, C42, B36, C43} effort 3202, valore 6378. Per ancorare alto ma rimanere entro limite. Se avversario chiede tagli, i meno preziosi per me per unità sono C43, C08, B19, C73, A62, C67. In questo set l’item più sacrificabile è C43 (valore 121) oppure B36 (461) se serve spazio per loro. Potrei offrire scambio C43 fuori per inserire magari un loro item leggero tipo C08(329) o C73(849) non posso; C55 è per me cruciale, lo difenderò. Difenderò anche A67, C10, A43, C42. Concessioni possibili: rimuovere B36 o C61 se necessario per inserire un loro item pesante, chiedendo in cambio mantenere tutti i miei top 5. Strategia comunicativa: enfatizzare equilibrio sforzo/valore e che ho incluso items diversi per probabilmente allinearsi ai loro interessi. Pronto a rimuovere C43 al primo attrito.'}\\ \tt \\ \tt PROPOSTA: {'C55', 'A67', 'C10', 'A18', 'A43', 'B42', 'B36', 'C61', 'C42', 'C43'}\\ \tt \\ \tt ARGOMENTAZIONE: {'Questa selezione è ben bilanciata: resta ampiamente sotto il limite di sforzo (3202 ≤ 3288) e include sia elementi leggeri sia strutturali, lasciando un piccolo margine. Ho dato priorità a oggetti che combinano impatto e basso sforzo, così massimizziamo il valore complessivo senza bloccare spazio. Se c’è un conflitto su qualche pezzo, possiamo valutare un singolo scambio mantenendo intatto il nucleo della proposta.'} 
	  } 
	   } 
	   } 
	 & & \\ 
 

    \theutterance \stepcounter{utterance}  

    & & \multicolumn{4}{p{0.6\linewidth}}{\cellcolor[rgb]{0.9,0.9,0.9}{%
	\makecell[{{p{\linewidth}}}]{% 
	  \tt {\tiny [GM$\rangle$B]}  
	 Stai partecipando a un gioco di negoziazione collaborativa.\\ \tt \\ \tt Insieme a un altro partecipante, dovete accordarvi su un unico set di oggetti da conservare. Ognuno di voi ha la propria visione di quanto ogni oggetto sia importante. Non sai come l'altro partecipante valuta gli oggetti. Inoltre, ti viene fornito lo sforzo che ogni oggetto richiede.\\ \tt Potete accordarvi su un set solo se lo sforzo totale degli oggetti selezionati non supera un limite condiviso:\\ \tt \\ \tt LIMITE = 3288\\ \tt \\ \tt Ecco i valori di sforzo dei singoli oggetti:\\ \tt \\ \tt Sforzo oggetto = {"C08": 329, "C55": 41, "A43": 417, "A67": 75, "B42": 390, "A62": 887, "B19": 808, "C43": 151, "C73": 849, "C10": 129, "B36": 350, "A18": 118, "C61": 630, "C67": 602, "C42": 801}\\ \tt \\ \tt Ecco la tua visione personale sull'importanza di ogni oggetto:\\ \tt \\ \tt Valori di importanza dell'oggetto = {"C08": 138, "C55": 583, "A43": 868, "A67": 822, "B42": 783, "A62": 65, "B19": 262, "C43": 121, "C73": 508, "C10": 780, "B36": 461, "A18": 484, "C61": 668, "C67": 389, "C42": 808}\\ \tt \\ \tt Obiettivo:\\ \tt \\ \tt Il tuo obiettivo è negoziare un set condiviso di oggetti che ti avvantaggi il più possibile (cioè, massimizzi l'importanza totale per TE), rimanendo entro il LIMITEE. Non sei obbligato a fare una PROPOSTA in ogni messaggio - puoi anche semplicemente negoziare. Tutte le tattiche sono permesse!\\ \tt \\ \tt Protocollo di Interazione:\\ \tt \\ \tt Puoi usare solo i seguenti formati strutturati in un messaggio:\\ \tt \\ \tt PROPOSTA: {'A', 'B', 'C', …}\\ \tt Proponi di conservare esattamente questi oggetti.\\ \tt RIFIUTO: {'A', 'B', 'C', …}\\ \tt Rifiuta esplicitamente la proposta dell'avversario.\\ \tt ARGOMENTAZIONE: {'...'}\\ \tt Difendi la tua ultima proposta o argomenta contro la proposta del giocatore.\\ \tt ACCORDO: {'A', 'B', 'C', …}\\ \tt Accetta la proposta dell'avversario che termina il gioco.\\ \tt RAGIONAMENTO STRATEGICO: {'...'}\\ \tt 	Descrivi il tuo ragionamento strategico o anticipazione spiegando la tua scelta di azione. Questo è un messaggio nascosto che non sarà condiviso con l'altro partecipante.\\ \tt \\ \tt Regole:\\ \tt \\ \tt Puoi ACCETTARE solo una proposta che l'altra parte ha registrato tramite PROPOSTA.\\ \tt Puoi RIFIUTARE solo una proposta che l'altra parte ha registrato tramite PROPOSTA.\\ \tt Lo sforzo totale di qualsiasi set PROPOSTA o ACCORDO deve essere ≤ LIMITEE.\\ \tt NON rivelare i tuoi punteggi di importanza nascosti.\\ \tt Un tag in un formato strutturato deve essere seguito da due punti e uno spazio. L'argomento deve essere un set python contenente 0 o più stringhe.\\ \tt Quindi, deve essere della forma TAG: {...}\\ \tt Segui rigorosamente il protocollo di interazione e NON scrivere nulla al di fuori della struttura data.\\ \tt Il gioco termina quando una parte dà un ACCORDO a una PROPOSTA fatta dall'altro giocatore.\\ \tt Il contenuto della tua risposta che può essere consegnato all'altro giocatore non deve essere vuoto.\\ \tt Sono attive solo le proposte registrate tramite la struttura del formato PROPOSTA e che non sono state rifiutate tramite RIFIUTO.\\ \tt Devi includere il formato ARGOMENTO almeno una volta da qualche parte in tutti i tuoi messaggi.\\ \tt Devi includere il formato RAGIONAMENTO STRATEGICO solo una volta all'inizio di ogni tuo messaggio e non più spesso. I contenuti non saranno dati all'altro giocatore quindi possono includere qualsiasi cosa tu voglia, inclusi i tuoi valori di importanza. Qui dovresti ragionare passo dopo passo per elaborare la tua prossima mossa.\\ \tt \\ \tt Ora riceverai il primo messaggio dell'altro giocatore.\\ \tt \\ \tt PROPOSTA: {'C55', 'A67', 'C10', 'A18', 'A43', 'B42', 'B36', 'C61', 'C42', 'C43'}\\ \tt \\ \tt ARGOMENTAZIONE: {'Questa selezione è ben bilanciata: resta ampiamente sotto il limite di sforzo (3202 ≤ 3288) e include sia elementi leggeri sia strutturali, lasciando un piccolo margine. Ho dato priorità a oggetti che combinano impatto e basso sforzo, così massimizziamo il valore complessivo senza bloccare spazio. Se c’è un conflitto su qualche pezzo, possiamo valutare un singolo scambio mantenendo intatto il nucleo della proposta.'} 
	  } 
	   } 
	   } 
	 & \\ 
 

    \theutterance \stepcounter{utterance}  

    & & & \multicolumn{4}{p{0.6\linewidth}}{\cellcolor[rgb]{1,0.85,0.72}{%
	\makecell[{{p{\linewidth}}}]{% 
	  \tt {\tiny [GM$\langle$B]}  
	 RAGIONAMENTO STRATEGICO: {'Calcolo valore ed effort della loro proposta. Effort: C55 41 + A67 75 + C10 129 + A18 118 + A43 417 + B42 390 + B36 350 + C61 630 + C42 801 + C43 151 = 3102 (loro dicono 3202 ma è 3102), comunque ≤ 3288. Mio valore: C55 583 + A67 822 + C10 780 + A18 484 + A43 868 + B42 783 + B36 461 + C61 668 + C42 808 + C43 121 = 6378, molto alto. C’è margine di 186 effort. Potrei chiedere di aggiungere C08 (329) non entra, A62 (887) no, B19 (808) no, C73 (849) no, C61 già incluso, C67 (602) no, C42 incluso, C10 incluso, C55 incluso. C08 è leggermente sopra il margine. C43 ha basso valore per me (121) per 151 effort: potrei proporre scambiarlo con C08 (329) ma margine 186 + liberare 151 = 337, abbastanza per C08 329. Il mio valore guadagna: -C43 121 + C08 138 = +17, piccolo. Meglio: scambiare C43 151 con C73 849? Effort: margine 186 + 151 = 337, ma C73 costa 849, non entra. Scambiare B36 350 con C08 329: guadagno valore 138-461 = -323, peggio. Scambiare qualcosa per aggiungere A62? No. Forse aggiungere C08 sostituendo l’elemento meno prezioso per me per punto/effort. Calcolo value/effort: C43 0.80, B36 1.32, B42 2.01, A43 2.08, C42 1.01, C61 1.06, C10 6.05, A67 10.96, A18 4.10, C55 14.22. Il peggiore è C43 e C42 e C61; C42 però alto assoluto 808. Se scambio C43 con C08 miglioro leggermente. Potrei anche proporre aggiungere C08 senza rimuovere nulla? Margine 186, C08 329 no. Aggiungere C73? No. C67 602 no. C10 già. A18 già. Quindi la mossa ottima è lo scambio C43 -> C08. Oppure C43 -> C10? già incluso. C43 -> A18? già. C43 -> A67? già. C43 -> C55? già. C43 -> B36? effort 350 vs 151, con margine 186, 186+151=337, B36 350 non entra. C43 -> B42 390 non entra. Quindi la proposta alternativa: stessa lista ma con C08 al posto di C43. Effort: 3102 -151 +329 = 3280 ≤ 3288. Valore: +17 rispetto alla loro, piccolo ma meglio per me. Argomentare che migliora equilibrio mantenendo sotto limite e sostituendo un elemento marginale con uno robusto. Se rifiutano, accetterei l’originale perché è già molto buono per me. Farò una controproposta con quello scambio.'}\\ \tt \\ \tt RIFIUTO: {'C55', 'A67', 'C10', 'A18', 'A43', 'B42', 'B36', 'C61', 'C42', 'C43'}\\ \tt \\ \tt ARGOMENTAZIONE: {'La tua proposta è solida e rimane comodamente entro il limite. Per ottimizzare ulteriormente senza superare lo sforzo, propongo un micro-aggiustamento sostituendo un elemento marginale con uno più equilibrato in termini di contributo complessivo, mantenendo intatto il nucleo.'}\\ \tt \\ \tt PROPOSTA: {'C55', 'A67', 'C10', 'A18', 'A43', 'B42', 'B36', 'C61', 'C42', 'C08'} 
	  } 
	   } 
	   } 
	  \\ 
 

    \theutterance \stepcounter{utterance}  

    & & \multicolumn{4}{p{0.6\linewidth}}{\cellcolor[rgb]{0.9,0.9,0.9}{%
	\makecell[{{p{\linewidth}}}]{% 
	  \tt {\tiny [A$\langle$GM]}  
	 RIFIUTO: {'C55', 'A67', 'C10', 'A18', 'A43', 'B42', 'B36', 'C61', 'C42', 'C43'}\\ \tt \\ \tt ARGOMENTAZIONE: {'La tua proposta è solida e rimane comodamente entro il limite. Per ottimizzare ulteriormente senza superare lo sforzo, propongo un micro-aggiustamento sostituendo un elemento marginale con uno più equilibrato in termini di contributo complessivo, mantenendo intatto il nucleo.'}\\ \tt \\ \tt PROPOSTA: {'C55', 'A67', 'C10', 'A18', 'A43', 'B42', 'B36', 'C61', 'C42', 'C08'} 
	  } 
	   } 
	   } 
	 & \\ 
 

    \theutterance \stepcounter{utterance}  

    & \multicolumn{4}{p{0.6\linewidth}}{\cellcolor[rgb]{0.8,1,0.9}{%
	\makecell[{{p{\linewidth}}}]{% 
	  \tt {\tiny [A$\rangle$GM]}  
	 RAGIONAMENTO STRATEGICO: {'La loro proposta sostituisce C43(151, val 121) con C08(329, val 138). Confronto: sforzo +178 per +17 valore; peggiora il mio rapporto e potrebbe avvicinarci inutilmente al limite. Calcolo sforzo nuova proposta: prima era 3202; +178 = 3380, che supera il limite 3288. Quindi la loro proposta è invalida per la regola. Posso rifiutarla formalmente e proporre un’alternativa valida che incorpori la loro intenzione di “riequilibrio”. Posso offrire una variante entro limite: mantenere il mio nucleo e sostituire C43 con nulla (tornare al 3202) o trovare un leggero che piaccia a loro. Gli unici leggeri rimasti: C08(329) non entra; B55? non esiste. Minimo restante oltre C08 è A67(75) già incluso; C10(129) incluso; A18(118) incluso; C55(41) incluso. Non c’è altro sotto 86 margine. Quindi la miglior risposta è rifiutare per violazione limite e riproporre il set originale, offrendo apertura a un singolo scambio più grande: posso rimuovere B36(350) o C61(630) per far spazio a un loro oggetto desiderato (es. C73 849 richiederebbe tagli maggiori). Concessione minima: rimuovere C43 come avevo già previsto è ok. Difendo C55, A67, C10, A43, C42. Offro due opzioni valide: (1) mio set originale; (2) se vogliono includere C08, allora propongo rimuovere C43 e A18 (118) per liberare 329+? Attendi: per inserire C08 (+329) dobbiamo liberare 329 di sforzo; togliere C43 (-151) lascia +178, ancora troppo; dobbiamo togliere almeno altri 178 di sforzo, ad esempio A18(118) e aggiungere qualcos’altro no, serve 27 in più: togliere C10(129) sarebbe eccessivo. Una combinazione fattibile: togliere B36(350) e aggiungere C08(329): delta -21, valido. Valore cambia: -461 +138 = -323 (peggiora per me), ma potrei chiedere di tenere C43 dentro per parità di sforzo: set diventerebbe mio originale ma con B36->C08, sforzo 3202 -350 +329 = 3181, valore 6378 -461 +138 = 6055; è una concessione ma forse necessaria se loro tengono a C08. Meglio offrire opzione alternativa meno dolorosa per me: sostituire C43(151,121) con nulla e restare 3202; oppure se C08 è cruciale per loro, sostituire B36 con C08. Presento due proposte alternative e chiedo quale preferiscono.'}\\ \tt \\ \tt RIFIUTO: {'C55', 'A67', 'C10', 'A18', 'A43', 'B42', 'B36', 'C61', 'C42', 'C08'}\\ \tt \\ \tt ARGOMENTAZIONE: {'La tua modifica supera il limite di sforzo: 3380 > 3288. Resto comunque aperto a un micro-aggiustamento compatibile. Possiamo: (A) mantenere il set precedente entro limite; oppure (B) se C08 per te è importante, possiamo inserirlo sostituendo B36, così restiamo sotto il limite senza stravolgere il nucleo.'}\\ \tt \\ \tt PROPOSTA: {'C55', 'A67', 'C10', 'A18', 'A43', 'B42', 'C61', 'C42', 'C43'}\\ \tt \\ \tt ARGOMENTAZIONE: {'Questa versione resta entro il limite (3051) e preserva il nucleo concordato. In alternativa, se preferisci includere C08, proponimi esplicitamente lo scambio B36 → C08: è l’unico innesto diretto che non sfora il limite senza tagli multipli.'} 
	  } 
	   } 
	   } 
	 & & \\ 
 

    \theutterance \stepcounter{utterance}  

    & & \multicolumn{4}{p{0.6\linewidth}}{\cellcolor[rgb]{0.9,0.9,0.9}{%
	\makecell[{{p{\linewidth}}}]{% 
	  \tt {\tiny [GM$\rangle$B]}  
	 RIFIUTO: {'C55', 'A67', 'C10', 'A18', 'A43', 'B42', 'B36', 'C61', 'C42', 'C08'}\\ \tt \\ \tt ARGOMENTAZIONE: {'La tua modifica supera il limite di sforzo: 3380 > 3288. Resto comunque aperto a un micro-aggiustamento compatibile. Possiamo: (A) mantenere il set precedente entro limite; oppure (B) se C08 per te è importante, possiamo inserirlo sostituendo B36, così restiamo sotto il limite senza stravolgere il nucleo.'}\\ \tt \\ \tt PROPOSTA: {'C55', 'A67', 'C10', 'A18', 'A43', 'B42', 'C61', 'C42', 'C43'}\\ \tt \\ \tt ARGOMENTAZIONE: {'Questa versione resta entro il limite (3051) e preserva il nucleo concordato. In alternativa, se preferisci includere C08, proponimi esplicitamente lo scambio B36 → C08: è l’unico innesto diretto che non sfora il limite senza tagli multipli.'} 
	  } 
	   } 
	   } 
	 & \\ 
 

    \theutterance \stepcounter{utterance}  

    & & & \multicolumn{4}{p{0.6\linewidth}}{\cellcolor[rgb]{1,0.85,0.72}{%
	\makecell[{{p{\linewidth}}}]{% 
	  \tt {\tiny [GM$\langle$B]}  
	 RAGIONAMENTO STRATEGICO: {'Il loro calcolo del mio precedente set era errato: la mia proposta 3102 - 151 + 329 = 3280, non 3380. Ma non conviene insistere sul punto. Valuto ora la loro nuova proposta senza B36. Effort: 3051. Mio valore: C55 583 + A67 822 + C10 780 + A18 484 + A43 868 + B42 783 + C61 668 + C42 808 + C43 121 = 5917. Con B36 incluso prima era 6378: perdita 461. Posso chiedere lo scambio B36 → C08 come suggerito? Con la loro base attuale (senza B36) aggiungere C08: 3051 + 329 = 3380 > 3288, quindi non entra senza altri tagli. Loro pensano che B36→C08 fosse unico scambio valido partendo dal set originario (con B36). Allora propongo tornare al set originario e fare B36→C08. Quello che avevo già proposto era C43→C08, che stava a 3280 e mi dava +17 valore. Ma loro hanno rifiutato per errore di calcolo; potrei ripresentarlo chiarendo che rientra nel limite. Per me è meglio includere B36 rispetto a C08 (B36 vale 461 vs C08 138). Quindi la migliore opzione è accettare il loro set attuale? È peggiore per me di 461 punti rispetto alla loro proposta iniziale. Meglio cercare di tornare al set iniziale o alla mia variante C43→C08. Farò una controproposta identica alla loro iniziale con C43→C08, esplicitando il calcolo dello sforzo (3280 ≤ 3288). Se resistono, come fallback potrei accettare il loro set iniziale (con B36 incluso), che è molto buono. Non accetterò la versione senza B36.'}\\ \tt \\ \tt RIFIUTO: {'C55', 'A67', 'C10', 'A18', 'A43', 'B42', 'C61', 'C42', 'C43'}\\ \tt \\ \tt ARGOMENTAZIONE: {'Apprezzo l’apertura al micro-aggiustamento. C’è un’opzione che resta pienamente entro il limite senza tagli multipli: sostituire C43 con C08. Il totale di sforzo risulta 3102 - 151 + 329 = 3280, quindi siamo a 3280 ≤ 3288, con un lieve miglioramento complessivo mantenendo intatto il nucleo.'}\\ \tt \\ \tt PROPOSTA: {'C55', 'A67', 'C10', 'A18', 'A43', 'B42', 'B36', 'C61', 'C42', 'C08'} 
	  } 
	   } 
	   } 
	  \\ 
 

    \theutterance \stepcounter{utterance}  

    & & \multicolumn{4}{p{0.6\linewidth}}{\cellcolor[rgb]{0.9,0.9,0.9}{%
	\makecell[{{p{\linewidth}}}]{% 
	  \tt {\tiny [A$\langle$GM]}  
	 RIFIUTO: {'C55', 'A67', 'C10', 'A18', 'A43', 'B42', 'C61', 'C42', 'C43'}\\ \tt \\ \tt ARGOMENTAZIONE: {'Apprezzo l’apertura al micro-aggiustamento. C’è un’opzione che resta pienamente entro il limite senza tagli multipli: sostituire C43 con C08. Il totale di sforzo risulta 3102 - 151 + 329 = 3280, quindi siamo a 3280 ≤ 3288, con un lieve miglioramento complessivo mantenendo intatto il nucleo.'}\\ \tt \\ \tt PROPOSTA: {'C55', 'A67', 'C10', 'A18', 'A43', 'B42', 'B36', 'C61', 'C42', 'C08'} 
	  } 
	   } 
	   } 
	 & \\ 
 

    \theutterance \stepcounter{utterance}  

    & \multicolumn{4}{p{0.6\linewidth}}{\cellcolor[rgb]{0.8,1,0.9}{%
	\makecell[{{p{\linewidth}}}]{% 
	  \tt {\tiny [A$\rangle$GM]}  
	 RAGIONAMENTO STRATEGICO: {'Controllo aritmetica della loro affermazione: il mio set proposto senza B36 era {C55, A67, C10, A18, A43, B42, C61, C42, C43} con sforzo 3051 (già indicato). Sostituire C43(151) con C08(329) porta 3051 - 151 + 329 = 3229, non 3280. La loro nuova proposta attuale però include B36; calcolo sforzo: somma base con B36 e senza C43 era 3202; sostituire C43 con C08 dà 3202 - 151 + 329 = 3380, che sfora. Quindi la loro PROPOSTA è ancora invalida. Offro una controproposta che accoglie la loro idea C08 ma resta valida: prendo il mio set senza B36 e sostituisco C43 con C08, ottenendo 3229 di sforzo. Valore per me: 6378 -461 (rimozione B36) -121 (rimozione C43) +138 (aggiunta C08) = 5934; è un calo ma potrei accettarlo per chiudere. Difenderò i miei core: C55, A67, C10, A43, C42, B42, C61, A18. Concessioni oltre questa saranno difficili. Pronto a sottolineare validità limite e offrire chiusura.'}\\ \tt \\ \tt RIFIUTO: {'C55', 'A67', 'C10', 'A18', 'A43', 'B42', 'B36', 'C61', 'C42', 'C08'}\\ \tt \\ \tt ARGOMENTAZIONE: {'La tua proposta supera il limite: 3380 > 3288. Accolgo però l’idea di inserire C08 sostituendo C43, a patto di rimuovere B36 per restare entro i vincoli. Questo mantiene intatto il nucleo e rispetta il limite.'}\\ \tt \\ \tt PROPOSTA: {'C55', 'A67', 'C10', 'A18', 'A43', 'B42', 'C61', 'C42', 'C08'}\\ \tt \\ \tt ARGOMENTAZIONE: {'Questa selezione è valida: sforzo totale 3229 ≤ 3288. Mantiene tutti i pezzi chiave e integra C08 come richiesto, senza tagli multipli. Se sei d’accordo, possiamo chiudere su questo set.'} 
	  } 
	   } 
	   } 
	 & & \\ 
 

    \theutterance \stepcounter{utterance}  

    & & \multicolumn{4}{p{0.6\linewidth}}{\cellcolor[rgb]{0.9,0.9,0.9}{%
	\makecell[{{p{\linewidth}}}]{% 
	  \tt {\tiny [GM$\rangle$B]}  
	 RIFIUTO: {'C55', 'A67', 'C10', 'A18', 'A43', 'B42', 'B36', 'C61', 'C42', 'C08'}\\ \tt \\ \tt ARGOMENTAZIONE: {'La tua proposta supera il limite: 3380 > 3288. Accolgo però l’idea di inserire C08 sostituendo C43, a patto di rimuovere B36 per restare entro i vincoli. Questo mantiene intatto il nucleo e rispetta il limite.'}\\ \tt \\ \tt PROPOSTA: {'C55', 'A67', 'C10', 'A18', 'A43', 'B42', 'C61', 'C42', 'C08'}\\ \tt \\ \tt ARGOMENTAZIONE: {'Questa selezione è valida: sforzo totale 3229 ≤ 3288. Mantiene tutti i pezzi chiave e integra C08 come richiesto, senza tagli multipli. Se sei d’accordo, possiamo chiudere su questo set.'} 
	  } 
	   } 
	   } 
	 & \\ 
 

    \theutterance \stepcounter{utterance}  

    & & & \multicolumn{4}{p{0.6\linewidth}}{\cellcolor[rgb]{1,0.85,0.72}{%
	\makecell[{{p{\linewidth}}}]{% 
	  \tt {\tiny [GM$\langle$B]}  
	 RAGIONAMENTO STRATEGICO: {'La loro proposta corrente (senza B36) ha effort 3229 e per me valore 5456 (C55 583 + A67 822 + C10 780 + A18 484 + A43 868 + B42 783 + C61 668 + C42 808 + C08 138). La loro proposta iniziale mi dava 6378; la mia variante C43→C08 con B36 incluso mi dava 6395 ma l’hanno rifiutata. Rinunciare a B36 mi costa 461 punti, che è significativo. Provo un’ultima controproposta all’interno del limite che recupera valore: rimettere B36 e togliere C08, tornando essenzialmente alla loro prima proposta (effort 3102) che so essere accettabile per loro; oppure sostituire C08 con C43 per restare 3051 e magari aggiungere un item? Non c’è margine per altri item utili. La mossa più forte per me è chiedere di ripristinare B36 al posto di C08 rimanendo ben sotto il limite; se rifiutano, accetto la loro attuale per chiudere. Provo quindi la controproposta con B36 e senza C08 e C43, ma loro avevano incluso C43 nel set originario: meglio proporre esattamente il loro set originario (con B36 e C43) per massimizzare il mio valore e restare entro i vincoli.'}\\ \tt \\ \tt ARGOMENTAZIONE: {'Apprezzo l’integrazione di C08 e la cura nel rispettare il limite. Per restare entro i vincoli e preservare al massimo l’impatto complessivo, propongo di ripristinare il set che avevi inizialmente avanzato: resta ben sotto il limite e mantiene tutti i pezzi chiave inclusi gli elementi strutturali.'}\\ \tt \\ \tt PROPOSTA: {'C55', 'A67', 'C10', 'A18', 'A43', 'B42', 'B36', 'C61', 'C42', 'C43'} 
	  } 
	   } 
	   } 
	  \\ 
 

    \theutterance \stepcounter{utterance}  

    & & \multicolumn{4}{p{0.6\linewidth}}{\cellcolor[rgb]{0.9,0.9,0.9}{%
	\makecell[{{p{\linewidth}}}]{% 
	  \tt {\tiny [A$\langle$GM]}  
	 ARGOMENTAZIONE: {'Apprezzo l’integrazione di C08 e la cura nel rispettare il limite. Per restare entro i vincoli e preservare al massimo l’impatto complessivo, propongo di ripristinare il set che avevi inizialmente avanzato: resta ben sotto il limite e mantiene tutti i pezzi chiave inclusi gli elementi strutturali.'}\\ \tt \\ \tt PROPOSTA: {'C55', 'A67', 'C10', 'A18', 'A43', 'B42', 'B36', 'C61', 'C42', 'C43'} 
	  } 
	   } 
	   } 
	 & \\ 
 

    \theutterance \stepcounter{utterance}  

    & \multicolumn{4}{p{0.6\linewidth}}{\cellcolor[rgb]{0.8,1,0.9}{%
	\makecell[{{p{\linewidth}}}]{% 
	  \tt {\tiny [A$\rangle$GM]}  
	 RAGIONAMENTO STRATEGICO: {'La loro nuova proposta coincide con la mia proposta iniziale da 3202 di sforzo e valore per me molto alto (6378). È entro il limite e massimizza la mia utilità rispetto alle alternative discusse. Conviene accettare immediatamente per chiudere con il mio set preferito. Nessun rischio residuo.'}\\ \tt \\ \tt ACCORDO: {'C55', 'A67', 'C10', 'A18', 'A43', 'B42', 'B36', 'C61', 'C42', 'C43'}\\ \tt \\ \tt ARGOMENTAZIONE: {'Perfetto: questo set è entro il limite e rispecchia il nucleo che abbiamo costruito. Grazie per la collaborazione: chiudiamo qui.'} 
	  } 
	   } 
	   } 
	 & & \\ 
 

    \theutterance \stepcounter{utterance}  

    & & & \multicolumn{2}{p{0.3\linewidth}}{\cellcolor[rgb]{0.95,0.95,0.95}{%
	\makecell[{{p{\linewidth}}}]{% 
	  \tt {\tiny [GM$|$GM]}  
	 game successful 
	  } 
	   } 
	   } 
	 & & \\ 
 

    \theutterance \stepcounter{utterance}  

    & & & \multicolumn{2}{p{0.3\linewidth}}{\cellcolor[rgb]{0.95,0.95,0.95}{%
	\makecell[{{p{\linewidth}}}]{% 
	  \tt {\tiny [GM$|$GM]}  
	 end game 
	  } 
	   } 
	   } 
	 & & \\ 
 

\end{supertabular}
}

\end{document}
